\begin{exercise}{5.1}{5}
    Man betrachte die komplexe Zahlenebene als euklidische Ebene und
    betrachte die Drehung $d$ mit Fixpunkt $p \in \mathbb C$ gegen den Uhrzeigersinn um den
    rechten Winkel, in Formeln gegeben durch $d : p + z\mapsto p + iz$. Man schreibe für
    einen beliebigen Richtungsvektor $w \in \mathbb C$ die Verknüpfung $(w+) \circ d$ wieder als
    eine Drehung. Was ist der Fixpunkt dieser neuen Drehung?
\end{exercise}

Zunächst schreiben dir $d$ für einzelne $x \in \mathbb C$ um, indem wir substituieren mit $x = p + z$, also $z = x -p $:
\[
    d(x) = p+i(x-p) = p+ix -ip= ix +p(1-i).
\]
$(w+)$ meint die Verschiebung, bzw. Translation um den Vektor $w$, also $x \mapsto x +w$.
Sei $d' = (w+) \circ d$, also ist $d'(x) = ix + p(1-i) +w$.
Um den neuen Fixpunkt $p'$ zu finden, lösen wir $ip'+p(1-i)+w=p'$:
\begin{align*}
    p' &= ip'+p(1-i)+w\\
    \Leftrightarrow p' - ip' &= p(1-i)+w\\
    \Leftrightarrow p'(1-i) &= p(1-i)+w\\
    \Leftrightarrow p' &= \frac{p(1-i)+w}{1-i} = p + \frac{w}{1-i}\\
    &= p + \frac{w(1+i)}{(1-i)(1+i)} = \boxed{p + \frac{1+i}{2}w}.
\end{align*}

\newpage
\begin{exercise}{5.3}{4}
    (\textbf{Dreispiegelungssatz}). Man zeige, dass sich jedes Element der Kongruenzgruppe einer euklidischen Ebene als Verknüpfung von einer, zwei oder drei
    Spiegelungen darstellen lässt.
\end{exercise}

\begin{minipage}{0.5\linewidth}
    Sei $f$ eine Isometrie (Kongruenz) einer euklidischen Ebene. Eine Isometrie wird eindeutig durch die Abbildung dreier nicht-kollinearen Punkte bestimmt.
    Seien also $A, B, C$ nicht-kollineare Punkte und $A' = f(A), B' = f(B'),$ und $C' = f(C')$ deren Bilder. Wir stellen nun $f$ als Verkettung dreier Spiegelungen dar:
    \begin{enumerate}
        \item Wenn $A \neq A'$, sei $\sigma_1$ die Reflektion an der Mittelsenkrechten von $AA'$. Wenn $A = A'$, sei $\sigma_1 = \operatorname{id}$. Also ist $\sigma_1(A) = A'$.
        \item Schreibe $B_1 := \sigma_1(B)$. Da Isometrien Abstände erhalten, gilt $A'B_1$ = $AB$ = $A'B'$. Somit liegt $A'$ auf der Mittelsenkrechten von $B'B_1$. Wenn $B_1 \neq B'$, dann sei $\sigma_2$ die Spiegelung an jener Mittelsenkrechte. Sonst sei $\sigma_2 = \operatorname{id}$. Somit gilt $\sigma_2(B_1) = B', \sigma_2(A') = A'$ und $(\sigma_2 \circ \sigma_1)$ bildet $A \mapsto A', B \mapsto B'$ ab. 
        \item Schreibe $C_2 := (\sigma_2 \circ \sigma_1)C$. Es gilt $A'C_2 = A'C'$ und $B'C_2 = B'C'$. Wenn $C_2 \neq C'$, ist $A'B'$ somit die Mittelsenkrechte von $C_2C'$. Analog zu vorher finden wir somit ein $\sigma_3$, welches $A'$, $B'$ fixiert und $\sigma_3(C_2) = C'$.
    \end{enumerate}
    Somit lässt sich $f$ schreiben als $\sigma_3 \circ \sigma_2 \circ \sigma_1$.
\end{minipage}
\begin{minipage}{0.49\linewidth}
    \begin{center}
        \fcolorbox{black}{white}{\begin{tikzpicture}[
        scale=1,
        point/.style={circle, fill, inner sep=1.5pt}
    ]

    % Punkte
    \coordinate (C) at (0,-0.2);
    \coordinate (B) at (1,1);
    \coordinate (A) at (1,-1);

    \coordinate (C') at (4,2.8);
    \coordinate (B') at (3,4);
    \coordinate (A') at (3,2);

    \coordinate (M) at (2,0.5);

    \coordinate (B1) at ({15/13},{16/13});
    \coordinate (C1) at ({122/65},{34/13});

    % Dreiecke
    \draw[thick, black!30] (B) -- (C) -- (A) -- cycle;
    \draw[thick, black!30] (B') -- (C') -- (A') -- cycle;

    % Sigma_1
    \draw[thick, red!40, -{Stealth}] (A) -- (A');
    \draw[dashed, thick, red]
    (-1,2.5) -- (5,-1.5)
    node[below right] {$\sigma_1$};

    % Punkte markieren
    \node[point, label=below left:$C$] at (C) {};
    \node[point, label=left:$B$] at (B) {};
    \node[point, label=below:$A$] at (A) {};
    \node[point, label=right:$C'$] at (C') {};
    \node[point, label=left:$B'$] at (B') {};
    \node[point, label=below right:$A'$] at (A') {};
            
\end{tikzpicture}}
        \fcolorbox{black}{white}{\begin{tikzpicture}[
        scale=1,
        point/.style={circle, fill, inner sep=1.5pt}
    ]

    % Punkte
    \coordinate (C) at (0,-0.2);
    \coordinate (B) at (1,1);
    \coordinate (A) at (1,-1);

    \coordinate (C') at (4,2.8);
    \coordinate (B') at (3,4);
    \coordinate (A') at (3,2);

    \coordinate (M) at (2,0.5);
    \coordinate (M2) at ({27/13},{34/13});

    \coordinate (B1) at ({15/13},{16/13});
    \coordinate (C1) at ({122/65},{34/13});
    \coordinate (C2) at (2, 2.8);

    % Dreiecke
    \draw[thick, black!30] (B) -- (C) -- (A) -- cycle;
    \draw[thick, black!30] (B') -- (C') -- (A') -- cycle;

    % Sigma_1
    \draw[dashed, thick, red!20]
    (-1,2.5) -- (5,-1.5)
    node[below right] {$\sigma_1$};

    % Sigma_2
    \draw[thick, blue, -{Stealth}] (B1) -- (B');
    \draw[dashed, thick, blue]
    ($(M2)-0.95*(3,-2)$) -- ($(M2)+0.95*(3,-2)$)
    node[right] {$\sigma_2$};

    % Punkte markieren
    \node[point, label=below left:$C$] at (C) {};
    \node[point, label=left:$B$] at (B) {};
    \node[point, label=below:$A$] at (A) {};
    \node[point, label=right:$C'$] at (C') {};
    \node[point, label=left:$B'$] at (B') {};
    \node[point, label=below right:$A'$] at (A') {};

    \node[point, black!40, label=right:$B_1$] at (B1) {};
      
\end{tikzpicture}}
        \fcolorbox{black}{white}{\begin{tikzpicture}[
        scale=1,
        point/.style={circle, fill, inner sep=1.5pt}
    ]

    % Punkte
    \coordinate (C) at (0,-0.2);
    \coordinate (B) at (1,1);
    \coordinate (A) at (1,-1);

    \coordinate (C') at (4,2.8);
    \coordinate (B') at (3,4);
    \coordinate (A') at (3,2);

    \coordinate (M) at (2,0.5);
    \coordinate (M2) at ({27/13},{34/13});

    \coordinate (B1) at ({15/13},{16/13});
    \coordinate (C1) at ({122/65},{34/13});
    \coordinate (C2) at (2, 2.8);

    % Dreiecke
    \draw[thick, black!30] (B) -- (C) -- (A) -- cycle;
    \draw[thick, black!30] (B') -- (C') -- (A') -- cycle;

    % Sigma_3
    \draw[thick, purple, -{Stealth}] (C2) -- (C');
    \draw[dashed, purple] ($(3, 5)$) -- ($(3,1)$);
    \node[label=below:\textcolor{purple}{$\sigma_3$}] at ($(3,1.2)$) {}; 
    \draw[dashed, thick, red!20]
    (-1,2.5) -- (5,-1.5)

    % Sigma_1
    node[below right] {$\sigma_1$};
    \draw[dashed, thick, blue!30]

    % Sigma_2
    ($(M2)-0.95*(3,-2)$) -- ($(M2)+0.95*(3,-2)$)
    node[right] {$\sigma_2$};

    % Punkte markieren
    \node[point, label=below left:$C$] at (C) {};
    \node[point, label=left:$B$] at (B) {};
    \node[point, label=below:$A$] at (A) {};
    \node[point, label=right:$C'$] at (C') {};
    \node[point, label=left:$B'$] at (B') {};
    \node[point, label=below right:$A'$] at (A') {};

    \node[point, label=below left:$C_1$] at (C1) {};
    \node[point, label=above right:$C_2$] at (C2) {};
    
\end{tikzpicture}}
    \end{center}
\end{minipage}

